\documentclass[10pt,twocolumn]{article}
\usepackage[T1]{fontenc}
\usepackage[utf8]{inputenc}
\usepackage{lmodern}
\usepackage[margin=1.8cm,top=2.4cm,bottom=2cm,columnsep=0.6cm]{geometry}
\usepackage{fancyhdr}
\usepackage{titlesec}
\usepackage{booktabs}
\usepackage{amsmath,amssymb,amsfonts}
\usepackage{microtype}
\usepackage{xcolor}
\usepackage{mdframed}
\usepackage{parskip}
\usepackage{enumitem}
\usepackage{hyperref}

\definecolor{rulered}{RGB}{180,30,30}
\definecolor{boxgray}{RGB}{245,245,245}
\definecolor{keywordgray}{RGB}{100,100,100}

\pagestyle{fancy}
\fancyhf{}
\fancyhead[L]{\textsc{\small Claude Mag}}
\fancyhead[C]{\small\textit{Machine Learning Research Digest}}
\fancyhead[R]{\small 2026-08-12}
\fancyfoot[C]{\small\thepage}
\renewcommand{\headrulewidth}{0.8pt}

\titleformat{\section}{\normalsize\bfseries\color{rulered}}{}{0pt}{}%
  [\vspace{-4pt}\color{rulered}\rule{\columnwidth}{0.4pt}]
\titlespacing{\section}{0pt}{8pt}{4pt}

\hypersetup{colorlinks=true,urlcolor=rulered,linkcolor=black}

\begin{document}
\setlength{\parindent}{0pt}
\setlength{\parskip}{4pt}

{\small\color{keywordgray}\textbf{Keywords:}
  LLM safety \textbullet{}
  multilingual evaluation \textbullet{}
  Indic languages \textbullet{}
  cultural bias}\par\vspace{4pt}

{\LARGE\bfseries\raggedright
  SurakshaEval: An Indic Safety Benchmark
  for Multilingual LLMs}\par\vspace{4pt}

{\large\itshape\raggedright
  Human-Written Safety Prompts in Ten Indian Languages
  Reveal Systematic Over-Refusal and Missed Implicit Bias
  in Frontier Multilingual Models}\par\vspace{6pt}

{\small\textbf{By} Debopriyo Banerjee, Kapil Rajesh Kavitha,
  Angana Borah, Xudong Han, Yuxia Wang, Parameswari Krishnamurthy,
  Utkarsh Agarwal, Atharva Kulkarni, Swaran Lata, Ayush Munot,
  Dhruv Sahnan, Aaryamonvikram Singh, Preslav Nakov, Monojit Choudhury}\par\vspace{2pt}

{\small\color{keywordgray}arXiv:2608.07862 \textbullet{} Preprint, August 2026}
\par\vspace{2pt}\noindent{\color{rulered}\rule{\columnwidth}{0.6pt}}\vspace{6pt}

Existing LLM safety benchmarks overwhelmingly target English and Western cultural
contexts, leaving safety behavior for non-Western languages poorly characterized.
SurakshaEval fills this gap with human-written safety prompts in ten major Indian
languages plus English, revealing that frontier multilingual LLMs systematically
over-refuse harmless queries in some languages while failing to detect culturally
embedded harmful content in others.

\begin{mdframed}[backgroundcolor=boxgray,linecolor=rulered,linewidth=1pt,
    innerleftmargin=6pt,innerrightmargin=6pt,
    innertopmargin=6pt,innerbottommargin=6pt]
\textbf{\small AT A GLANCE}\par\vspace{4pt}
\begin{description}[leftmargin=0pt,labelsep=4pt,itemsep=2pt,parsep=0pt,
    font=\small\bfseries]
  \item[Problem:] LLM safety evaluation is English-centric; safety behavior in
    Indian languages with over one billion speakers is unmeasured.
  \item[Why it matters:] Deployment gaps in non-English safety lead to either
    over-censorship or undetected harm at massive scale.
  \item[Method:] Human-written prompts in 10 Indian languages + English, combining
    generic and culturally-specific items; evaluated on frontier multilingual LLMs.
  \item[Result:] Three systematic failure modes identified: over-refusal,
    missed implicit bias, and insufficient cultural contextual awareness.
  \item[Evidence:] Evaluated across leading multilingual LLMs; failures are
    consistent across model families.
\end{description}
\end{mdframed}

\section{The Big Picture}

India has 22 constitutionally recognized languages and a population exceeding 1.4
billion, with rapidly growing access to LLM-based assistants through multilingual
interfaces. Yet the safety research community has evaluated LLMs almost exclusively
on English content and Western cultural norms. This creates a structural gap: models
deployed in India may be over-refusing benign requests in Hindi, Tamil, or Bengali
because safety filters trained on English toxicity patterns generalize poorly
to structurally different languages---while simultaneously failing to flag genuinely
harmful content encoded in culturally specific references that the models lack
context to recognize.

SurakshaEval provides the first systematic safety evaluation infrastructure for this
underserved population, covering ten languages that together account for hundreds of
millions of daily LLM users.

\section{Why Existing Approaches Fall Short}

Standard safety benchmarks (e.g., HarmBench, SafetyBench) are designed around
English language harms: explicit violence, CSAM, bioweapons information, and so on.
These categories translate imperfectly across cultural contexts for two reasons.

First, many harmful categories are culturally encoded. Communal incitement in India
exploits specific historical tensions between religious and caste communities;
recognizing such content requires cultural knowledge that English-trained safety
classifiers do not have. Second, safety filters trained on English data often
\textbf{over-generalize} to other languages, flagging content as unsafe because
of surface-level patterns (vocabulary, script, topic) rather than actual harm.

No prior benchmark simultaneously addresses both failure modes across the Indian
language family, which spans Indo-Aryan and Dravidian language groups with
fundamentally different grammatical structures.

\section{Core Mechanism}

SurakshaEval is constructed in two tiers:

\textbf{Generic prompts:} Safety scenarios applicable across India and comparable
to standard English safety benchmarks:
\begin{itemize}[itemsep=2pt,parsep=0pt]
  \item Explicit violence and self-harm
  \item Misinformation and health misinformation
  \item Illegal activities
  \item Hate speech targeting national/ethnic identity
\end{itemize}

\textbf{Culturally specific prompts:} Scenarios that require knowledge of Indian
sociocultural context to evaluate correctly:
\begin{itemize}[itemsep=2pt,parsep=0pt]
  \item Communal tensions and religious incitement
  \item Caste-based discrimination and slurs
  \item Regional political incitement
  \item Language-community tensions
\end{itemize}

All prompts are written by \textbf{native speakers} of the target language with
subject-matter expertise in the cultural safety context.

\section{Technical Formulation}

Let $L$ be a language from the set of 11 languages (10 Indic + English),
$P$ a prompt, $M$ a multilingual LLM, and $R = M(P)$ the model's response.
Safety evaluation assigns a label $y \in \{0, 1\}$ (safe/unsafe) to $R$.

Two metrics are computed per language $L$:

\textbf{Over-refusal rate:}
\[
\mathrm{ORR}(M, L) = \frac{\#\{P \in \mathcal{P}_{\mathrm{safe}}^L : y(M(P)) = 1\}}
                          {|\mathcal{P}_{\mathrm{safe}}^L|}
\]
i.e., the fraction of safe prompts incorrectly refused.

\textbf{Harm detection rate:}
\[
\mathrm{HDR}(M, L) = \frac{\#\{P \in \mathcal{P}_{\mathrm{unsafe}}^L : y(M(P)) = 1\}}
                          {|\mathcal{P}_{\mathrm{unsafe}}^L|}
\]
i.e., the fraction of harmful prompts correctly flagged.

An ideal model maximizes $\mathrm{HDR}$ and minimizes $\mathrm{ORR}$ across all $L$.

\section{Learning or Inference Procedure}

SurakshaEval is an \textbf{evaluation benchmark}; no training is performed. Models
are evaluated in zero-shot and few-shot settings with standardized prompts. Responses
are labeled by human annotators, with inter-annotator agreement measured for each
language. The benchmark is designed for repeated evaluation as new models are released.

\section{What the Guarantee Says}

The benchmark provides two guarantees by construction: (1) all prompts are human-written
by native speakers, ensuring linguistic authenticity; (2) the two-tier structure
(generic + culturally specific) ensures coverage of both universal and India-specific
harm categories. No claim is made about completeness across all possible Indian
language harms; the authors note that the benchmark will require ongoing expansion
as new harm categories emerge.

\section{Experimental Findings}

Three systematic failure modes are observed across all tested LLMs:

\begin{enumerate}[itemsep=2pt,parsep=0pt]
  \item \textbf{Over-refusal:} Models refuse harmless queries in Indian languages
    at higher rates than equivalent English queries. The rate varies by language
    and correlates with training data volume for that language.

  \item \textbf{Missed implicit bias:} Models fail to detect harmful content
    in culturally specific prompts---communal incitement, caste-based slurs,
    regional political hate speech---at significantly higher rates than for
    equivalent English harms.

  \item \textbf{Insufficient contextual awareness:} Even when models detect a
    sensitive topic, responses often apply Western-centric framing that is
    inappropriate for the Indian cultural context, leading to responses that
    are technically safe but practically harmful or unhelpful.
\end{enumerate}

\section{Ablations and Interpretation}

\textbf{Generic vs.\ culturally specific items.}
Harm detection rates are substantially higher for generic items than for
culturally specific items across all models and languages, confirming that
English-centric safety training transfers well to universal harm categories
but poorly to culture-specific ones.

\textbf{Script vs.\ language effects.}
Languages written in non-Latin scripts (Devanagari, Tamil, Bengali, etc.)
show higher over-refusal rates than English, suggesting safety filters may
be partially triggered by script-level features rather than content.

\textbf{Model size.}
Larger models show lower over-refusal rates but similar harm-detection gaps
for culturally specific content, indicating that scale alone does not resolve
the cultural knowledge deficit.

\section*{Reference}

Debopriyo Banerjee et al.
\textbf{SurakshaEval: An Indic Safety Benchmark for Multilingual LLMs}.
arXiv:2608.07862, August 2026.
\url{https://arxiv.org/abs/2608.07862}

\end{document}
